\phantomsection
\addcontentsline{toc}{chapter}{Abstract}
\markboth{Abstract}{Abstract}

\begin{center}
\textbf{\Large Abstract}
\end{center}

Retrieval-Augmented Generation (RAG) gives Large Language Models (LLMs) access
to external knowledge, yet multi-hop questions remain challenging: they require
several dependent pieces of evidence, and not every question demands the same
amount of retrieval and reasoning. This thesis investigates how such systems
can become more adaptive and evidence-aware, from recognising compositional
difficulty before answering to improving the retrieval process once question
decomposition is used.

First, we ask whether an LLM's internal representations already reveal how
compositionally difficult a question is. Across three language models and three
multi-hop benchmarks, benchmark-defined difficulty is linearly decodable from
hidden representations, emerges progressively across model depth, survives
controls for label permutation and selected surface features, exhibits ordinal
structure, and transfers more consistently across benchmarks after surface
residualisation. These findings identify a potential pre-answer signal for
adaptive computation.

The thesis then turns to question decomposition. We introduce an evidence-aware
evaluation framework that separates final-answer correctness from whether the
required supporting evidence is retrieved, how deeply it appears in the
rankings, and which subquestions are responsible for retrieving it. The
evidence-ownership estimator recovers 92 of 100 complete reference-style
assignment masks. On a controlled MuSiQue development panel, decomposition
increases evidence availability and raises GPT-4o-mini exact match from 40\% to
52\%, while further experiments show that better retrieval does not always
translate directly into better answers.

Finally, the proposed measurements are used to study automated decomposition
prompt search and adaptive passage allocation. Prompt optimisation reveals that
retrieval-oriented and answer-oriented objectives can favour different
instructions. Under a tight eight-passage budget, learned allocation improves
complete-evidence coverage by 6.3 percentage points on 2WikiMultiHopQA and
4.7 points on MuSiQue, while a separate 700-question MuSiQue experiment yields
a 2.1-point F1 gain. Together, these results support a view of multi-hop RAG in
which reasoning strategy, decomposition quality, and retrieval resources are
adapted to the structure and evidence needs of each question.

\vspace{0.8em}

\noindent\textbf{Keywords:} Retrieval-Augmented Generation, Large Language
Models, Multi-Hop Question Answering, Question Decomposition, Compositional
Difficulty, Representation Probing, Adaptive Retrieval, Evidence-Aware
Evaluation.